\subsection{LGS mit nicht-invertierbarer Matrix}

Im vorangegangenen Abschnitt haben wir gesehen, dass man Gleichungssysteme leicht mit Hilfe der inversen Matrix lösen kann. In diesem Kapitel wollen wir den Fall untersuchen, dass die Matrix nicht invertierbar ist. 

\begin{LA}
{}{}
In dieser Lernaufgabe untersuchen wir Gleichungssysteme der Form
$$ A\cdot \begin{pmatrix}
    x\\y
\end{pmatrix}= \vec{b} \qquad \text{ bzw. } \qquad B\cdot \begin{pmatrix}
    x\\y
\end{pmatrix}= \vec{b}$$
mit den (nicht invertierbaren) Matrizen
$$ A = \begin{pmatrix}
    4 & -6\\
    -6 & 9
\end{pmatrix} \qquad \text{ bzw. } \qquad B = \begin{pmatrix}
    12 & 20\\
    9 & 15
\end{pmatrix}$$


\textbf{Teil 1}: der Kern der Matrix

\begin{enumerate}
    \item Finden Sie drei Lösungen der Gleichung $A\cdot \begin{pmatrix}
    x\\y
\end{pmatrix}= \begin{pmatrix}
    0\\0
\end{pmatrix}$
\item Zeigen Sie, dass die drei Lösungen alle auf einer Geraden liegen und diese Gerade der Menge $\left\{t\cdot \begin{pmatrix}
    3\\2
\end{pmatrix} \mid t \in \R\right\}$ entspricht.
\item Betrachten Sie nun die Matrix $B$. Bestimmen die Lösungsmenge der Gleichung $B\cdot \begin{pmatrix}
    x\\y
\end{pmatrix}= \begin{pmatrix}
    0\\0
\end{pmatrix}$.
\end{enumerate}

\textbf{Teil 2:} das Bild der Matrix
\begin{enumerate}
    \item Berechnen Sie $A \cdot \vec{v}$ für drei verschiedene Vektoren $\vec{v} \in \R^2$.
    \item Zeigen Sie, dass die drei Vektoren $A \cdot \vec{v}$ alle auf einer Geraden liegen und diese Gerade der Menge $\left\{t\cdot \begin{pmatrix}
    2\\-3
\end{pmatrix} \mid t \in \R\right\}$ entspricht.
    \item Betrachten Sie nun die Matrix $B$. Bestimmen Sie die Menge aller Bilder $B \cdot \vec{v}$ für $\vec{v} \in \R^2$.

\end{enumerate}

\textbf{Teil 3:}
\begin{enumerate}
    \item Erklären Sie in Bezugnahme auf `Teil 2', wieso das Gleichungssystem 
    $$A\cdot \begin{pmatrix}
    x\\y
\end{pmatrix}= \begin{pmatrix}
    10\\10
\end{pmatrix}$$
keine Lösung hat.
\item Der Vektor $\begin{pmatrix}
    5\\-2
\end{pmatrix}$ ist eine Lösung des Gleichungssystems
$$A\cdot \begin{pmatrix}
    x\\y
\end{pmatrix}= \begin{pmatrix}
    32\\-48
\end{pmatrix}$$
Erklären Sie in Bezugnahme auf `Teil 1', wieso jeder Vektor der Form
$$\begin{pmatrix}
    5\\-2
\end{pmatrix} + t \cdot \begin{pmatrix}
    3\\2
\end{pmatrix}$$
mit $t \in \R$ eine Lösung des Gleichungssystems ist.
\end{enumerate}

\textbf{Teil 4:}

Entscheiden Sie, ob die folgenden Gleichungssysteme Lösungen haben und wenn ja, geben Sie die Lösungsmenge an.
\begin{enumerate}
    \item $B \cdot \begin{pmatrix}
    x\\y
\end{pmatrix}= \begin{pmatrix}
    6\\8
\end{pmatrix}$
    \item $B \cdot \begin{pmatrix}
    x\\y
\end{pmatrix}= \begin{pmatrix}
    8\\6
\end{pmatrix}$
\end{enumerate}
\end{LA}





















\clearpage
Im folgenden werden einige neue Begriffe eingeführt, mit denen man die Lösungsmenge eines LGS mit nicht-invertierbarer Matrix beschreiben kann. Diese Begriffe sind wichtig, um die Lösungsmenge eines LGS zu verstehen und zu beschreiben.

\begin{definition}
    {Kern und Bild}{}
   
Sei \( A\) eine $2\times2$-Matrix. 

Die Menge aller Vektoren, die auf den Nullvektor abgebildet werden, nennt man den \emph{Kern} von \( A \) und schreibt dafür:
\[
\ker(A) = \left\{ \vec{v} \in \mathbb{R}^2 \mid A\cdot \vec{v} = \vec{0} \right\}
\]

Die Menge aller Bilder $A\cdot \vec{v}$ von \( A \) nennt man das \emph{Bild} von \( A \) und schreibt dafür:
\[
\operatorname{im}(f) = \left\{ A\cdot \vec{v} \mid \vec{v} \in \mathbb{R}^2 \right\}
\]
Die Bezeichnung \( \operatorname{im} \) steht für \emph{image} (englisch für Bild).


\end{definition}


\begin{bsp}
    {Kern und Bild von $A = \begin{pmatrix} 1 & 3 \\ 3 & 9 \end{pmatrix}$ }{}


Dreimal die erste Spalte minus die zweite ergibt den Nullvektor, also ist:
\[
\ker(A) = \left\{ t \begin{pmatrix} 3 \\ -1 \end{pmatrix} \;\middle|\; t \in \mathbb{R} \right\}.
\]

Beide Spalten sind Vielfache des Vektors $\begin{pmatrix} 1 \\ 3 \end{pmatrix}$.  
Also gilt:
\[
\operatorname{im}(A) = \left\{ t \begin{pmatrix} 1 \\ 3 \end{pmatrix} \;\middle|\; t \in \mathbb{R} \right\}.
\]

\end{bsp}

\begin{bsp}
{Kern und Bild von $B = \begin{pmatrix} 1 & 3 \\ 4 & -1 \end{pmatrix}.$}{}

Die Matrix ist invertierbar, weil $\det(B)  = -13 \neq 0$. 
Also gilt:

\[
\ker(B) = \left\{B^{-1}\cdot \vec{0} \right\} = \left\{ \vec{0} \right\}.
\]
und 
\[
\operatorname{im}(B)  = \mathbb{R}^2.
\]

\end{bsp}

\begin{bsp}
{Kern und Bild von $N = \begin{pmatrix} 0 & 0 \\ 0 & 0 \end{pmatrix}.$}{}

Die Matrix $N$ bildet jeden Vektor auf den Nullvektor ab.
Also gilt:
\[
\ker(N) = \R^2.
\]
und 
\[
\operatorname{im}(N)  = \left\{\vec{0}\right\}.
\]

\end{bsp}

Mit diesen Begriffen kann man die Lösungsmenge eines LGS mit nicht-invertierbarer Matrix beschreiben.
\begin{satz}{Lösung von LGS }{}
Sei \( A \) eine $2\times2$-Matrix und \( \vec{b} \in \mathbb{R}^2 \). Dann gilt:
\begin{itemize}
    \item Falls $A$ invertierbar ist, dann hat das LGS \( A\cdot \vec{x} = \vec{b} \) genau eine Lösung, nämlich \( \vec{x} = A^{-1}\cdot \vec{b} \).
    \item Falls $A$ nicht invertierbar ist, dann gilt:
\begin{itemize}
    \item Wenn \( \vec{b} \notin \operatorname{im}(A) \), dann hat das LGS \( A\cdot \vec{x} = \vec{b} \) keine Lösung.
    \item Wenn \( \vec{b} \in \operatorname{im}(A) \), dann hat das LGS \( A\cdot \vec{x} = \vec{b} \) unendlich viele Lösungen. 
    \item Ist \( \vec{x}_0 \) eine Lösung des LGS \( A\cdot \vec{x} = \vec{b} \), dann ist die Menge aller Lösungen gegeben durch:
    \[
    \left\{ \vec{x}_0 + \vec{v} \mid \vec{v} \in \ker(A) \right\}.
    \]
\end{itemize}
\end{itemize}
\end{satz}

\begin{bsp}
{ $ \begin{pmatrix} 1 & 3 \\ 3 & 9 \end{pmatrix} \cdot \vec{x} = \begin{pmatrix}
    8\\ 16
\end{pmatrix}$ }{}

\vspace*{7 cm}
% Sei $A = \begin{pmatrix} 1 & 3 \\ 3 & 9 \end{pmatrix} $. 
% Das Bild von $A$ ist die Gerade durch den Nullvektor in Richtung $\begin{pmatrix} 1 \\ 3 \end{pmatrix}$.
% Der Vektor $\begin{pmatrix} 8 \\ 16 \end{pmatrix}$ liegt nicht auf dieser Geraden, also hat das LGS keine Lösung.
\end{bsp}

\begin{bsp}
{$\begin{pmatrix} 1 & 3 \\ 3 & 9 \end{pmatrix} \cdot \vec{x} = \begin{pmatrix}
    8\\ 24
\end{pmatrix}$ }{}

\vspace*{12 cm}

% Sei wieder $A = \begin{pmatrix} 1 & 3 \\ 3 & 9 \end{pmatrix} $. 
% Der Vektor $\begin{pmatrix} 8 \\ 24 \end{pmatrix}$ liegt im Bild von $A$. 

% Der Kern von $A$ ist gegeben durch:
% \[
% \ker(A) = \left\{ t \begin{pmatrix} 3 \\ -1 \end{pmatrix} \;\middle|\; t \in \mathbb{R} \right\}.
% \]


% Eine Lösung des LGS ist $\begin{pmatrix} 2 \\ 2 \end{pmatrix}$.

% Die gesamte Lösungsmenge ist gegeben durch:
% \[
% \left\{ \begin{pmatrix} 2 \\ 2 \end{pmatrix} + t \begin{pmatrix} 3 \\ -1 \end{pmatrix} \;\middle|\; t \in \mathbb{R} \right\}.
% \]
\end{bsp}

\clearpage

\begin{ueb}
    {}{}

\begin{enumerate}
  \item Bestimmen Sie den Kern und das Bild der folgenden Matrizen.  
  \begin{enumerate}
    \item \( A = \begin{pmatrix} 1 & 2 \\ 3 & 4 \end{pmatrix} \)
    \item \( B = \begin{pmatrix} 0 & 0 \\ 0 & 0 \end{pmatrix} \)
    \item \( C = \begin{pmatrix} 2 & 12 \\ -1 & -6 \end{pmatrix} \)
    \item \( D = \begin{pmatrix} -6 & 7 \\ 12 & -14 \end{pmatrix} \)
    \item \( E = \begin{pmatrix} 1 & 3 \\ 0 & 0 \end{pmatrix} \)
  \end{enumerate}

  \item Geben Sie die Lösungsmenge der folgenden LGS an. Die Matrizen beziehen sich auf die Aufgabe 1.
  \begin{enumerate}
    \item \( A \cdot \begin{pmatrix} x \\ y \end{pmatrix} = \begin{pmatrix} 1 \\ 2 \end{pmatrix} \)
    \item  \( B \cdot \begin{pmatrix} x \\ y \end{pmatrix} = \begin{pmatrix} 0 \\ 0 \end{pmatrix} \)
    \item \( C \cdot \begin{pmatrix} x \\ y \end{pmatrix} = \begin{pmatrix} 3 \\ 20 \end{pmatrix} \)
    \item \( D \cdot \begin{pmatrix} x \\ y \end{pmatrix} = \begin{pmatrix} -5 \\ 10 \end{pmatrix} \)
    \item \( E \cdot \begin{pmatrix} x \\ y \end{pmatrix} = \begin{pmatrix} 1 \\ 0 \end{pmatrix} \)
  \end{enumerate}

  \item Sei \( f \) die Projektion auf die Gerade \( y = x \) in Richtung  
  \( \vec{v} = \begin{pmatrix} 1 \\ -1 \end{pmatrix} \).  
  Bestimmen Sie den Kern und das Bild von \( f \).

  \item Geben Sie eine Matrix an, deren Kern die Gerade  
  \( g \colon \vec{r}(t) = t \begin{pmatrix} 1 \\ 2 \end{pmatrix} \)  
  und deren Bild die \( y \)-Achse ist.

  \item Geben Sie eine Matrix \( A \) an, deren Kern und Bild übereinstimmen.  
  Was ergibt \( A^2 \)?

  \item Sei \( f \) die Spiegelung an der \( xz \)-Ebene.  
  Bestimmen Sie den Kern und das Bild.

  \item Sei \( f \) die orthogonale Projektion auf die \( xz \)-Ebene.  
  Bestimmen Sie den Kern und das Bild.

  \item \textbf{Challenge:}
  
    Beweisen Sie die folgende Aussage:  
    Sei \( f \) eine lineare Abbildung mit \( \ker(f) = \{ \vec{0} \} \).  
    Dann ist \( f \) injektiv.\\
    \emph{Tipp:} Injektiv bedeutet: Aus \( f(\vec{v}) = f(\vec{w}) \) folgt \( \vec{v} = \vec{w} \).  
    Betrachten Sie \( f(\vec{v} - \vec{w}) \).

\end{enumerate}

\end{ueb}

\clearpage
\emph{Lösungen:}
\begin{enumerate}
  \item 
  \begin{enumerate}
    \item \( \ker(A) = \left\{ \vec{0} \right\} \), \( \operatorname{im}(A) = \mathbb{R}^2 \)
    \item \( \ker(B) = \mathbb{R}^2 \), \( \operatorname{im}(B) = \left\{ \vec{0} \right\} \)
    \item \( \ker(C) = \left\{ t\begin{pmatrix} 6 \\ -1 \end{pmatrix} | t\in\mathbb{R}  \right\}  ,  \operatorname{im}(C) =  \left\{ t\begin{pmatrix} 2 \\ -1  \end{pmatrix} | t\in\mathbb{R}  \right\}\)
    \item \( \ker(D) = \left\{ t\begin{pmatrix} 7 \\ 6 \end{pmatrix} | t\in\mathbb{R}  \right\}  ,  \operatorname{im}(D) =  \left\{ t\begin{pmatrix} -1 \\ 2  \end{pmatrix} | t\in\mathbb{R}  \right\}\)
    \item \( \ker(E) = \left\{ t\begin{pmatrix} 3 \\ -1 \end{pmatrix} | t\in\mathbb{R}  \right\}  ,  \operatorname{im}(E) =  \left\{ t\begin{pmatrix} 1 \\ 0  \end{pmatrix} | t\in\mathbb{R}  \right\}\)
    \end{enumerate}
    \item \begin{enumerate}
        \item $\LL = \left\{\begin{pmatrix}
            0 \\ \frac{1}{2}
        \end{pmatrix}\right\}$
        \item $\LL = \R^2$
        \item $\LL = \emptyset$
        \item $\LL = \left\{\begin{pmatrix}
            \frac{5}{6} \\ 0
        \end{pmatrix} + t\begin{pmatrix}
            7 \\ 6
        \end{pmatrix} \mid t \in \R\right\}$
        \item $\LL = \left\{\begin{pmatrix}
            1 \\ 0
        \end{pmatrix} + t\begin{pmatrix}
            3 \\ -1
        \end{pmatrix} \mid t \in \R\right\}$
    \end{enumerate}
    \item \( \ker(f) = \left\{ t \begin{pmatrix} 1 \\ -1 \end{pmatrix} \;\middle|\; t \in \mathbb{R} \right\} \) und $\operatorname{im}(f) = \left\{ t \begin{pmatrix} 1 \\ 1 \end{pmatrix} \;\middle|\; t \in \mathbb{R} \right\}$.
    \item Z.B. \( \begin{pmatrix} 0 & 0 \\ 2 & -1 \end{pmatrix} \)
    \item Z.B. \( \begin{pmatrix} 0 & 0 \\  1 & 0\end{pmatrix} \). Dann ist logischerweise $A^2 = \begin{pmatrix} 0 & 0 \\  0 & 0\end{pmatrix}$.
    \item \( \ker(f) = \left\{ \vec{0} \right\} \) und \( \operatorname{im}(f) = \R^3 \)
    \item \( \ker(f) = \left\{ \begin{pmatrix} 0 \\ y \\ 0 \end{pmatrix} \;\middle|\; y \in \mathbb{R} \right\} \) und \( \operatorname{im}(f) = \left\{ \begin{pmatrix} x \\ 0 \\ z \end{pmatrix} \;\middle|\; x,z \in \mathbb{R} \right\} \)
\end{enumerate}